\documentclass[tikz,border=12pt]{standalone}
\usepackage{fontspec}
\usepackage{xeCJK}
\IfFontExistsTF{Segoe UI}{\setmainfont{Segoe UI}}{\setmainfont{Arial}}
\IfFontExistsTF{Microsoft JhengHei}{\setCJKmainfont{Microsoft JhengHei}}{\setCJKmainfont{Noto Sans CJK TC}}
\xeCJKsetup{CJKglue={\hskip 0.09em plus 0.02em}, CJKecglue={\hskip 0.12em plus 0.02em}}
\usetikzlibrary{arrows.meta,calc,positioning}
\definecolor{paper}{HTML}{FFFDF8}
\definecolor{ink}{HTML}{0F172A}
\definecolor{muted}{HTML}{475569}
\definecolor{line}{HTML}{CBD5E1}
\definecolor{blue}{HTML}{2563EB}
\definecolor{bluebg}{HTML}{EAF2FF}
\definecolor{green}{HTML}{00856A}
\definecolor{greenbg}{HTML}{E6FAF2}
\definecolor{gold}{HTML}{D97706}
\definecolor{goldbg}{HTML}{FFF2C7}
\definecolor{red}{HTML}{DC2626}
\definecolor{redbg}{HTML}{FFE8E8}
\tikzset{
  title/.style={font=\fontsize{16.5}{22}\selectfont\bfseries,text=ink,align=center},
  sub/.style={font=\fontsize{9.8}{14}\selectfont,text=muted,align=center},
  item/.style={rounded corners=13pt,line width=1pt,align=center,inner xsep=11pt,inner ysep=8pt,font=\fontsize{9.8}{13.4}\selectfont,text=ink},
  label/.style={font=\fontsize{11.5}{15}\selectfont\bfseries,text=ink,align=center},
  small/.style={font=\fontsize{8.8}{12}\selectfont,text=muted,align=center}
}
\begin{document}
\begin{tikzpicture}[x=1cm,y=1cm,line cap=round,line join=round]
\path[fill=paper] (0,0) rectangle (16,9);

\node[title,text width=13.8cm] at (8,8.16) {邊界畫清楚，內部才有空間學};
\node[sub,text width=12.2cm] at (8,7.5) {自我改進不是放任，而是把可改與不可自改分開。};

\draw[rounded corners=22pt,draw=red,line width=1.15pt,fill=redbg] (0.95,1.45) rectangle (15.05,6.65);
\node[label] at (2.45,5.9) {外層守門};
\node[item,fill=white,draw=red,text width=2.5cm] at (6.15,5.68) {權限};
\node[item,fill=white,draw=red,text width=2.5cm] at (10.65,5.68) {正式評分};
\node[item,fill=white,draw=red,text width=2.5cm] at (6.15,4.72) {保留測試};
\node[item,fill=white,draw=red,text width=2.5cm] at (10.65,4.72) {稽核紀錄};

\draw[rounded corners=20pt,draw=green,line width=1.15pt,fill=greenbg] (3.15,2.05) rectangle (12.85,3.92);
\node[label] at (4.42,3.43) {內層學習};
\node[item,fill=bluebg,draw=blue,text width=1.55cm] at (5.95,2.72) {提示詞};
\node[item,fill=goldbg,draw=gold,text width=1.55cm] at (8.55,2.72) {工具順序};
\node[item,fill=white,draw=line,text width=1.55cm] at (11.15,2.72) {錯誤處理};

\node[small,text width=10.8cm] at (8,1.02) {外層不是裝飾，它決定 agent 不能用哪種方式讓自己看起來進步。};
\end{tikzpicture}
\end{document}
